\documentclass[fontsize=11pt]{article}

\usepackage[margin=1in]{geometry}
\usepackage{graphicx}
\usepackage{wrapfig}
\usepackage{subcaption}
\usepackage{amsmath} 
\usepackage{siunitx}
\usepackage{booktabs}
\usepackage{gensymb}
\usepackage{hyperref}
\usepackage[export]{adjustbox}
\usepackage{anyfontsize}

\begin{document}

	\graphicspath{{./figures/}}

	%BJC: check address, phone number and fax (?) number
	\begin{figure}
	\centering
	\begin{minipage}{0.37\textwidth}
	\includegraphics[width=2.14in,left]{CUBoulder.pdf}
	\end{minipage}
	\begin{minipage}{0.62\textwidth}
	\scriptsize
	\noindent Department of Chemical and Biological Engineering \hfill t 303-735-7860~~~~~~~~~~~~~~~~~~ \\
	\noindent 596 UCB \hfill f 303-492-8425~~~~~~~~~~~~~~~~~~ \\
	\noindent Boulder, Colorado 80309 \hfill michael.shirts@colorado.edu \\
	\end{minipage}
	\end{figure}
	
	%\noindent May 3\textsuperscript{rd}, 2019\\
	\noindent \today \\

	\noindent Dear Editors,\\
	
	\newcommand{\ManuscriptTitle}{Capturing Subdiffusive Solute Dynamics and 
	Predicting Selectivity in Nanoscale Pores with Time Series Modeling}
	
	We are submitting ``\ManuscriptTitle'' by Benjamin J. Coscia and Michael R.
	Shirts, for consideration for publication as an Article in the Journal of 
	Physical Chemistry B. We expect this work to be of significant interest to 
	researchers who study complex transport in both structured and amorphous 
	soft materials. Please note that this article has been uploaded to arXiv with
	differences only in formatting: \href{https://arxiv.org/abs/2004.07905}{https://arxiv.org/abs/2004.07905}. \\
	
	In this article, we use solute trajectories generated from 5$\mu$s molecular
	dynamics (MD) simulations of an inverted hexagonal (H\textsubscript{II}) phase
	self-assembled lyotropic liquid crystal (LLC) polymer membrane in order to 
	parameterize two different stochastic time series models capable of predicting
	solute flux and selectivity on a macroscopic scale. Realizations of both models
	quantitatively reproduce the mean squared displacements (MSDs) exhibited by 
	solutes in MD simulations. This is among the first studies to apply time 
	series modeling of molecular simulations in order to predict bulk materials properties.\\

	This article builds on recent work of ours also published in J. Phys. Chem. B 
    where we characterized transport mechanisms exhibited by a series of 
	small polar solutes in an LLC Membrane (Coscia and Shirts \textit{J. Phys.
	Chem. B}, 123, 6314--6330 (2019)). In that study, we observed subdiffusive solute transport
	characterized by intermittent hops separated by periods of entrapment. In
	this work, we build on the chemical intuition gained from our previous study
	in order to formulate stochastic models which produce trajectory realizations
	that share the same dynamical properties as those exhibited by solutes in 
	our MD simulations. We explore two modeling approaches. First, based on the
	anomalous diffusion literature, we parameterize solute motion as 
    subordinated fractional Brownian and L\'evy motion processes, based on the observation that 
	solutes experience a series of dwell times separated by anti-correlated hop
	lengths. Our second approach adds state-dependent dynamics to standard Markov
	state modeling approaches. We model solute motion as a series of transitions
	between various trapped states, each with different dynamical properties.
	Finally, we demonstrate how one can project our best model onto macroscopic 
	timescales in order to estimate solute flux and selectivity.\\
	
	%BJC: included phone numbers because it asked to do so in author's guide.
	\noindent Some suggestions for reviewers are:
	\begin{enumerate}
	
		\item Francisco Hung has used molecular simulations to study nematic liquid crystals
		and transport under confinement within various nanostructured materials 
		(Northeastern University, 617-373-8619,\\ \href{mailto:f.hung@northeastern.edu}{f.hung@northeastern.edu})
		
		\item Mahesh Mahanthappa has experience modeling soft materials using molecular dynamics
		including work on ion transport in various ordered liquid crystalline phases. 
		(University of Minnesota, 612-625-4599, \href{mailto:maheshkm@umn.edu}{maheshkm@umn.edu})
		
		\item Stephen Paddison uses various levels of molecular simulation theory to understand
		proton transport in nanostructured and charged polymer membranes. (The University of Tennessee 
		Knoxville, 865-974-2026, \href{mailto:spaddison@utk.edu}{spaddison@utk.edu})
		
	\end{enumerate}
	
	\noindent Please send correspondence regarding this paper to Michael R. Shirts (contact
	details in letterhead).\\	
	
	\noindent Sincerely,
	
	\noindent Michael R. Shirts \\
	\noindent Benjamin J. Coscia \\
	
\end{document}

% LocalWords:  BJC Subdiffusive Solute Nanoscale Coscia arXiv solute lyotropic
% LocalWords:  LLC MSDs solutes Phys Chem subdiffusive evy timescales Klafter
% LocalWords:  univeristy JPCB woould nanostructured Mahesh Mahanthappa klafter
% LocalWords:  Tel Paddison
